\section[Stage 236]{Stage 236: Physical calibration and material thresholds}
\label{stage:236}

\stagefield{Verification}{SymPy audit: \StageFile{scripts/moving_throat_pde_stage236_physical_calibration_and_material_threshold_companion_from_the_stage235_export_and_cold_survival_compiler_sympy_audit.py}.  Mathematica audit: \StageFile{mathematica/moving_throat_pde_stage236_physical_calibration_and_material_threshold_companion_from_the_stage235_export_and_cold_survival_compiler_mathematica_audit.wl}.}

\subsection{Physical calibration dictionary}
\label{subsec:app-part08-stage236-calibration}

Stage 236 introduces the explicit physical dictionary
\begin{equation}
\label{eq:app-part08-stage236-physical-dictionary}
t^{\rm phys}=t_*t,
\qquad
r^{\rm phys}=\frac{\lambda_{\rm phys}}{\lambda_{\rm ref}}r,
\qquad
E^{\rm phys}=E_*E.
\end{equation}
The safe-edge lattice event-equivalent rate is
\begin{equation}
\label{eq:app-part08-stage236-lat-safe-rate}
\gamma_{\rm lat,safe}^{\rm eq}
=
f_{\rm lat}(s_c)\mu_\eta\frac{s_0^2-s_c^2}{s_c}.
\end{equation}
A coarse transport projection factor \(\Upsilon_{\rm lat}>0\) defines the physical lattice-turnover rate
\begin{equation}
\label{eq:app-part08-stage236-lattice-turn-rate}
\gamma_{\rm lat,turn}^{\rm phys}
=
\frac{\gamma_{\rm lat,safe}^{\rm eq}}{\Upsilon_{\rm lat}t_*}
=
\zeta_{\rm ep}\lambda_{\rm ep}\omega_D.
\end{equation}
Thus the electron-phonon turnover threshold is \cref{eq:app-part08-main-ep-threshold}, and the product form is
\begin{equation}
\label{eq:app-part08-stage236-ep-product}
(\lambda_{\rm ep}\omega_Dt_{\rm cross}^{\rm phys})_{\min}^{(236)}
=
\frac{f_{\rm lat}(s_c)\mu_\eta(s_0^2-s_c^2)}{\Upsilon_{\rm lat}\zeta_{\rm ep}s_c^2}.
\end{equation}
The legacy Session-V slice is recovered by choosing
\begin{equation}
\label{eq:app-part08-stage236-upsilon-session}
\Upsilon_{\rm lat}^{\rm sess}
=
\frac{\gamma_{\rm lat,safe}^{\rm eq}}{\gamma_{\rm lattice}^{\rm red}}.
\end{equation}
This is a calibration slice, not a new branch law.

For the harmonic interstitial companion,
\begin{equation}
\label{eq:app-part08-stage236-physical-radius}
r_{\rm turn}^{\rm phys}=\frac{r_{\rm turn}}{\lambda_{\rm ref}}\lambda_{\rm phys},
\end{equation}
and
\begin{equation}
\label{eq:app-part08-stage236-chi-lattice}
\chi_{\lambda,{\rm lattice}}(r_{\rm turn})
=
\frac{2\lambda_{\rm phys}}{r_{\rm turn}^{\rm phys}}
=
\frac{2\lambda_{\rm ref}}{r_{\rm turn}}.
\end{equation}
This is only a geometric trigger.  The force-matched stiffness requires
\begin{equation}
\label{eq:app-part08-stage236-force-match}
k_{\rm eff,req}
=
\frac{E_*\lambda_{\rm ref}^{2}}{\lambda_{\rm phys}^{2}}
\frac{|V_{\rm red}'(r_{\rm turn})|}{r_{\rm turn}}
=
\mathcal K_{\rm turn}\frac{E_*}{\lambda_{\rm phys}^{2}},
\end{equation}
where
\begin{equation}
\label{eq:app-part08-stage236-kturn}
\mathcal K_{\rm turn}=\frac{\lambda_{\rm ref}^{2}}{r_{\rm turn}}|V_{\rm red}'(r_{\rm turn})|.
\end{equation}
For \(\lambda_{\rm phys}=a_{\rm int}/2\),
\begin{equation}
\label{eq:app-part08-stage236-k-aint}
k_{\rm eff,req}=4\mathcal K_{\rm turn}\frac{E_*}{a_{\rm int}^{2}}.
\end{equation}

The Korringa companion is the cleanest physical map.  With
\begin{equation}
\label{eq:app-part08-stage236-korringa}
T_1T=\mathcal K_{\rm corr},
\qquad
T_1\ge t_{\rm cross}^{\rm phys},
\end{equation}
the ceiling temperature is
\begin{equation}
\label{eq:app-part08-stage236-tmax}
T_{\max}=\frac{\mathcal K_{\rm corr}}{t_{\rm cross}^{\rm phys}}
=\frac{s_c\mathcal K_{\rm corr}}{t_*}.
\end{equation}

\subsection{Material-screening ratios}
\label{subsec:app-part08-stage236-screening-ratios}

The companion ends with four screening ratios:
\begin{equation}
\label{eq:app-part08-stage236-pi-ep}
\Pi_{\rm ep}
=
\frac{\Upsilon_{\rm lat}\zeta_{\rm ep}\lambda_{\rm ep}\omega_Dt_*}{\gamma_{\rm lat,safe}^{\rm eq}}
=
\frac{\lambda_{\rm ep}\omega_D}{(\lambda_{\rm ep}\omega_D)_{\min}^{(236)}},
\end{equation}
\begin{equation}
\label{eq:app-part08-stage236-pi-chi}
\Pi_\chi=\frac{2\lambda_{\rm ref}}{r_{\rm turn}},
\end{equation}
\begin{equation}
\label{eq:app-part08-stage236-pi-k}
\Pi_k=
\frac{k_{\rm eff}\lambda_{\rm phys}^{2}}{\mathcal K_{\rm turn}E_*}
=
\frac{k_{\rm eff}}{k_{\rm eff,req}},
\end{equation}
\begin{equation}
\label{eq:app-part08-stage236-pi-t}
\Pi_T(T)=
\frac{\mathcal K_{\rm corr}}{Tt_{\rm cross}^{\rm phys}}
=
\frac{T_{\max}}{T}.
\end{equation}
The materials companion is then
\begin{equation}
\label{eq:app-part08-stage236-final-screen}
\Pi_{\rm ep}\ge1,
\qquad
\Pi_\chi\ge1,
\qquad
\Pi_k\ge1,
\qquad
\Pi_T(T)\ge1.
\end{equation}
This is not a claim that any particular host passes.  It is the reduced threshold stack that a candidate host must satisfy after the physical unit map and material constants are supplied.
